%----------------------------------------------------------------------------------------
% NOTATIE EN FIGUREN VAN DE CURSUS GELIJKVORMIGHEID (2de graad)
% Staat in de preambule en niet in een hoofdstuk: ook de lijsten met oplossingen gebruiken dit.
% Vertrekpunt: notatie.tex van Kegelsneden, zonder de krommen en kegels.
%----------------------------------------------------------------------------------------
% Bewaarde grafiek, geplaatst door maak-pdf.py: bestand en pagina in figuren-cache/
\newcommand{\grafiekbestand}[2]{\includegraphics[page=#2,trim=3bp 3bp 3bp 3bp]{figuren-cache/#1}}
\usetikzlibrary{angles,quotes}

% Lengte van een lijnstuk: \lengte{AB} geeft |AB|; \abs{-2} de absolute waarde (\Abs met grotere strepen).
% Met \lvert en \rvert, zodat een minteken na de streep geen aftrekking wordt.
\newcommand{\lengte}[1]{\lvert #1\rvert}
\newcommand{\abs}[1]{\lvert #1\rvert}
\newcommand{\Abs}[1]{\bigl\lvert #1\bigr\rvert}
% Driehoek en hoek: \drieh{ABC} geeft de driehoek ABC, \hoek{A} een hoek met een dakje
\newcommand{\drieh}[1]{\triangle #1}
\newcommand{\hoek}[1]{\widehat{#1}}
% Gelijkvormig: \gv geeft het teken ~
\newcommand{\gv}{\sim}
% Schaal 1 : 50 met een smalle spatie: \schaal{1}{50}
\newcommand{\schaal}[2]{#1\,{:}\,#2}

% Zoals stappen, maar met een bredere kolom voor de berekening: voor lange formules
\newenvironment{stappenbreed}{%
  \par\renewcommand{\arraystretch}{1.25}%
  \noindent\begin{tabular}{@{}>{$\displaystyle}p{9.8cm}<{$}@{\hspace{4mm}}>{\RaggedRight\small\itshape\arraybackslash}p{4.9cm}@{}}%
}{\end{tabular}\par}

% Zoals stappen, met eigen kolombreedtes: \begin{stappenvrij}{6.2cm}{3.8cm}. Voor naast een figuur.
\newenvironment{stappenvrij}[2]{%
  \par\renewcommand{\arraystretch}{1.25}%
  \noindent\begin{tabular}{@{}>{$\displaystyle}p{#1}<{$}@{\hspace{3mm}}>{\RaggedRight\small\itshape\arraybackslash}p{#2}@{}}%
}{\end{tabular}\par}

% Tekst of berekening links, figuur rechts: \naast{links}{rechts}
\newsavebox{\naastbox}
\newcommand{\naast}[3][]{%
  \par\noindent\sbox{\naastbox}{#3}%
  \begin{minipage}[t]{\dimexpr\linewidth-\wd\naastbox-3mm\relax}\vspace{0pt}#2\end{minipage}\hfill
  \begin{minipage}[t]{\wd\naastbox}\vspace{0pt}\usebox{\naastbox}\end{minipage}\par\addvspace{1mm}}
% Een rij figuren, gecentreerd
\newcommand{\figuren}[1]{\par\addvspace{1mm}{\centering #1\par}\addvspace{1.5mm}}
% Opgave met een figuur rechts: \begin{metgrafiek}[breedte]{tikz} opgave en deelvragen \end{metgrafiek}
\newsavebox{\grafiekbox}
\newenvironment{metgrafiek}[2][]{%
  \sbox{\grafiekbox}{#2}%
  \begin{minipage}[t]{\dimexpr\linewidth-\wd\grafiekbox-3mm\relax}%
}{%
  \end{minipage}\hfill
  \begin{minipage}[t]{\wd\grafiekbox}\vspace{-2mm}\usebox{\grafiekbox}\end{minipage}\par}

% Rooster zonder assen, voor schaaltekeningen: \ruitjes{xmin}{ymin}{xmax}{ymax}
\newcommand{\ruitjes}[4]{\draw[black!14,very thin] (#1,#2) grid (#3,#4);}
% Een punt als stip: \stip{x}{y}, of met label: \punt[positie]{x}{y}{label}
\newcommand{\stip}[2]{\fill (#1,#2) circle (1.5pt);}
\newcommand{\punt}[4][above]{\fill (#2,#3) circle (1.5pt) node[#1,inner sep=2pt] {$#4$};}
% Een hoekboogje bij hoekpunt B tussen de stralen naar A en C (tegenwijzerzin van A naar C):
% \hoekboog[opties]{A}{B}{C}{straal}; werkt met benoemde coördinaten
\newcommand{\hoekboog}[5][]{\pic[draw,accent,angle radius=#5,#1] {angle=#2--#3--#4};}
% Een rechte hoek bij B tussen A en C: \rechtehoek{A}{B}{C}
\newcommand{\rechtehoek}[3]{\pic[draw,grijs,angle radius=2.2mm] {right angle=#1--#2--#3};}
% Een maatstok voor een kaart: \maatstok{x}{y}{lengte in cm}{onderverdeling}{label rechts}
\newcommand{\maatstok}[5]{%
  \foreach \i [evaluate=\i as \xa using #1+(\i-1)*#3/#4, evaluate=\i as \xb using #1+\i*#3/#4] in {1,...,#4}{%
    \pgfmathtruncatemacro{\even}{mod(\i,2)}%
    \ifnum\even=1 \fill[grijs] (\xa,#2) rectangle (\xb,#2+.12);\else\draw[grijs] (\xa,#2) rectangle (\xb,#2+.12);\fi}
  \draw[grijs] (#1,#2) rectangle (#1+#3,#2+.12);
  \node[right,font=\scriptsize] at (#1+#3,#2+.06) {#5};}

% Ruimtelijke figuren in schuin aanzicht: optie [schuin] bij tikzpicture. De z-as wijst naar achter (schuin omhoog).
\tikzset{schuin/.style={x={(1cm,0cm)},y={(0cm,1cm)},z={(.42cm,.3cm)}}}
% Een kubusje met ribbe 1 en de hoek vooraan linksonder in (x, y, z): \blokje{x}{y}{z}.
% Teken de achterste eerst (grote z), daarna van onder naar boven en van links naar rechts.
\newcommand{\blokje}[3]{%
  \filldraw[fill=accent!22,draw=accent,line join=round] (#1,#2,#3) -- ++(1,0,0) -- ++(0,1,0) -- ++(-1,0,0) -- cycle;
  \filldraw[fill=accent!10,draw=accent,line join=round] (#1,#2+1,#3) -- ++(1,0,0) -- ++(0,0,1) -- ++(-1,0,0) -- cycle;
  \filldraw[fill=accent!40,draw=accent,line join=round] (#1+1,#2,#3) -- ++(0,0,1) -- ++(0,1,0) -- ++(0,0,-1) -- cycle;}
% Een balk met afmetingen a (breed), b (hoog), c (diep), hoek vooraan linksonder in (x, y): \balk{x}{y}{a}{b}{c}
\newcommand{\balk}[5]{%
  \filldraw[fill=accent!22,draw=accent,line join=round] (#1,#2,0) -- ++(#3,0,0) -- ++(0,#4,0) -- ++(-#3,0,0) -- cycle;
  \filldraw[fill=accent!10,draw=accent,line join=round] (#1,#2+#4,0) -- ++(#3,0,0) -- ++(0,0,#5) -- ++(-#3,0,0) -- cycle;
  \filldraw[fill=accent!40,draw=accent,line join=round] (#1+#3,#2,0) -- ++(0,0,#5) -- ++(0,#4,0) -- ++(0,0,-#5) -- cycle;}
% Een paashaas met hoogte 12 (in de eenheden van de figuur), voeten in (x, y): \haas{x}{y}{schaal}
\newcommand{\haas}[3]{%
  \begin{scope}[shift={(#1,#2)},scale=#3]
    \filldraw[fill=accent!22,draw=accent,thick] (-.2,10.9) ellipse [x radius=.6,y radius=1.25,rotate=12];
    \filldraw[fill=accent!22,draw=accent,thick] (1.4,10.8) ellipse [x radius=.6,y radius=1.25,rotate=-18];
    \filldraw[fill=accent!22,draw=accent,thick] (0,3.6) ellipse [x radius=3,y radius=3.6];
    \filldraw[fill=accent!22,draw=accent,thick] (.6,8.2) circle (2);
    \fill[accent] (1.3,8.6) circle (.22);
    \filldraw[fill=accent!22,draw=accent,thick] (-3,2.2) circle (1);
  \end{scope}}
